\makeglossaries

% =========================================================
% Metodología
% =========================================================
\newacronym{DSRM}{DSRM}{\textit{Design Science Research Methodology}}
\newacronym{AHP}{AHP}{\textit{Analytic Hierarchy Process}}
% \newacronym{DPI}{DPI}{\textit{Inspección profunda de paquetes}}



% =========================================================
% Núcleo OT / ICS
% =========================================================
\newacronym{ICS}{ICS}{\textit{Industrial Control System}}
\newacronym{IACS}{IACS}{\textit{Industrial Automation and Control System}}
\newacronym{OT}{OT}{\textit{Operational Technology}}
\newacronym{IT}{IT}{\textit{Information Technology}}
\newacronym{SCADA}{SCADA}{\textit{Supervisory Control and Data Acquisition}}
\newacronym{DCS}{DCS}{\textit{Distributed Control System}}
\newacronym{PLC}{PLC}{\textit{Programmable Logic Controller}}
\newacronym{RTU}{RTU}{\textit{Remote Terminal Unit}}
\newacronym{IED}{IED}{\textit{Intelligent Electronic Device}}
\newacronym{HMI}{HMI}{\textit{Human-Machine Interface}}
\newacronym[description={Controlador de estación de la Micro-Red; en este trabajo no debe interpretarse como \textit{Application Programming Interface}}]{API}{API}{\textit{Advanced Power Integrated Station Controller}}
\newacronym{SBC}{SBC}{\textit{Single Board Computer}}
\newacronym{SoC}{SoC}{\textit{System-on-Chip}}
\newacronym{BITW}{BITW}{\textit{Bump-in-the-Wire}}
\newacronym{HITL}{HITL}{\textit{Hardware-in-the-Loop}}
\newacronym{IoT}{IoT}{\textit{Internet of Things}}
\newacronym{IIoT}{IIoT}{\textit{Industrial Internet of Things}}
\newacronym{DT}{DT}{\textit{Digital Twin}}
\newglossaryentry{AGRO5}{name={AGRO5}, description={Laboratorio y proyecto de simulación de cadena agroalimentaria orientado a Industria 5.0}}


% =========================================================
% Redes, direccionamiento y transporte
% =========================================================
\newacronym{OSI}{OSI}{\textit{Open Systems Interconnection}}
\newacronym{TCP}{TCP}{\textit{Transmission Control Protocol}}
\newacronym{IP}{IP}{\textit{Internet Protocol}}
\newacronym{IPv4}{IPv4}{\textit{Internet Protocol version 4}}
\newacronym{LAN}{LAN}{\textit{Local Area Network}}
\newacronym{WAN}{WAN}{\textit{Wide Area Network}}
\newacronym{DMZ}{DMZ}{\textit{Demilitarized Zone}}
\newacronym{NAT}{NAT}{\textit{Network Address Translation}}
\newacronym{AP}{AP}{\textit{Access Point}}
\newglossaryentry{ethernet}{name={Ethernet}, description={Tecnología de red de capa de enlace utilizada para transportar tramas en redes locales}}
\newglossaryentry{bridge}{name={bridge}, description={Puente de capa 2 que reenvía tramas entre interfaces de red sin operar como router}}
\newacronym{MTU}{MTU}{\textit{Maximum Transmission Unit}}
\newacronym{PMTU}{PMTU}{\textit{Path MTU}}
\newacronym{MSS}{MSS}{\textit{Maximum Segment Size}}
\newacronym{ICMP}{ICMP}{\textit{Internet Control Message Protocol}}
\newacronym{PCAP}{PCAP}{\textit{Packet capture}}
\newacronym{ARP}{ARP}{\textit{Address Resolution Protocol}}
\newacronym{TTL}{TTL}{\textit{Time To Live}}
\newacronym{TDM}{TDM}{\textit{Time Division Multiplexing}}
\newglossaryentry{RS232}{name={RS-232}, description={Estándar de comunicación serial usado para enlaces punto a punto entre dispositivos}}
\newglossaryentry{RS485}{name={RS-485}, description={Estándar de comunicación serial diferencial usado en buses multipunto industriales}}


% =========================================================
% Protocolos industriales y Modbus
% =========================================================
\newacronym{MODBUS}{Modbus}{\textit{Modbus Protocol}}
\newacronym{MODBUSTCP}{Modbus TCP}{\textit{Modbus over TCP/IP}}
\newglossaryentry{MODBUSRTU}{name={Modbus RTU}, description={Variante serial de Modbus que usa una representación binaria compacta de los mensajes}}
\newglossaryentry{MODBUSASCII}{name={Modbus ASCII}, description={Variante serial de Modbus que codifica los mensajes mediante caracteres ASCII}}
\newglossaryentry{modbusSecurityProfile}{name={Modbus/TCP Security}, description={Perfil seguro para Modbus TCP que emplea TLS, certificados X.509 y autenticación mutua}}
\newacronym{MBAP}{MBAP}{\textit{Modbus Application Protocol Header}}
\newacronym{ADU}{ADU}{\textit{Application Data Unit}}
\newacronym{PDU}{PDU}{\textit{Protocol Data Unit}}
\newacronym{ASCII}{ASCII}{\textit{American Standard Code for Information Interchange}}
\newacronym{OPCUA}{OPC UA}{\textit{Open Platform Communications Unified Architecture}}
\newacronym{DNP3SA}{DNP3-SA}{\textit{DNP3 Secure Authentication}}
\newacronym{RPC}{RPC}{\textit{Remote Procedure Call}}
\newacronym{HTTPS}{HTTPS}{\textit{Hypertext Transfer Protocol Secure}}
\newacronym{MQTTS}{MQTTS}{\textit{Message Queuing Telemetry Transport Secure}}


% =========================================================
% Túneles, canales seguros y seguridad de red
% =========================================================
\newacronym{VPN}{VPN}{\textit{Virtual Private Network}}
\newacronym{TLS}{TLS}{\textit{Transport Layer Security}}
\newacronym{SSH}{SSH}{\textit{Secure Shell}}
\newacronym{IPsec}{IPsec}{\textit{Internet Protocol Security}}
\newacronym{ESP}{ESP}{\textit{Encapsulating Security Payload}}
\newacronym{IKE}{IKE}{\textit{Internet Key Exchange}}
\newacronym{MACsec}{MACsec}{\textit{Media Access Control Security}}
\newacronym{MKA}{MKA}{\textit{MACsec Key Agreement}}
\newglossaryentry{OpenVPN}{name={OpenVPN}, description={Software de red privada virtual que puede operar como túnel de capa 3 o puente de capa 2}}
\newglossaryentry{Wireguard}{name={WireGuard}, description={Protocolo y software de VPN de capa 3 orientado a configuración simple y bajo sobrecosto}}

% =========================================================
% Criptografía, llaves y certificados
% =========================================================
\newacronym{CIA}{CIA}{\textit{Confidentiality, Integrity and Availability}}
\newacronym{AES}{AES}{\textit{Advanced Encryption Standard}}
\newacronym{RSA}{RSA}{\textit{Rivest-Shamir-Adleman}}
\newacronym{ECC}{ECC}{\textit{Elliptic Curve Cryptography}}
\newacronym{HMAC}{HMAC}{\textit{Hash-Based Message Authentication Code}}
\newacronym{SHA}{SHA}{\textit{Secure Hash Algorithm}}
\newacronym{AEAD}{AEAD}{\textit{Authenticated Encryption with Associated Data}}
\newacronym{PKI}{PKI}{\textit{Public Key Infrastructure}}
\newacronym{CA}{CA}{\textit{Certificate Authority}}
\newacronym{CRL}{CRL}{\textit{Certificate Revocation List}}
\newacronym{X509}{X.509}{\textit{X.509 Public Key Certificate}}
\newacronym{PSK}{PSK}{\textit{Pre-Shared Key}}
\newglossaryentry{cifradoSimetrico}{name={cifrado simétrico}, description={Mecanismo criptográfico que emplea una clave compartida para cifrar y descifrar información}}
\newglossaryentry{criptografiaClavePublica}{name={criptografía de clave pública}, description={Modelo criptográfico basado en pares de claves pública y privada para autenticación, intercambio de claves o firma digital}}
\newglossaryentry{funcionHash}{name={función hash criptográfica}, plural={funciones hash criptográficas}, description={Función que genera un resumen de longitud fija a partir de datos de entrada para apoyar verificaciones de integridad}}
\newglossaryentry{autenticacionMutua}{name={autenticación mutua}, description={Verificación criptográfica recíproca de identidad entre ambos extremos de una comunicación}}

% =========================================================
% Amenazas y monitoreo de seguridad
% =========================================================
\newacronym{MITM}{MITM}{\textit{Man-in-the-Middle}}
\newacronym{DoS}{DoS}{\textit{Denial of Service}}
\newacronym{DPI}{DPI}{\textit{Deep Packet Inspection}}
\newacronym{eBPF}{eBPF}{\textit{extended Berkeley Packet Filter}}
\newglossaryentry{sniffing}{name={sniffing}, description={Captura pasiva de tráfico de red para observar metadatos o contenido cuando no está protegido}}
\newglossaryentry{spoofing}{name={spoofing}, description={Suplantación de información de red para inducir una asociación o identidad falsa}}
\newglossaryentry{replay}{name={replay}, description={Reinyección de tráfico previamente válido fuera de su contexto temporal original}}
\newacronym{IDS}{IDS}{\textit{Intrusion Detection System}}
\newacronym{IPS}{IPS}{\textit{Intrusion Prevention System}}
\newacronym{RBAC}{RBAC}{\textit{Role-Based Access Control}}
\newacronym{DDoS}{DDoS}{\textit{Distributed Denial of Service}}
\newacronym{CVE}{CVE}{\textit{Common Vulnerabilities and Exposures}}
\newacronym{NVT}{NVT}{\textit{Network Vulnerability Test}}

% =========================================================
% Métricas de desempeño
% =========================================================
\newacronym{RTT}{RTT}{\textit{Round-Trip Time}}
\newacronym{CPU}{CPU}{\textit{Central Processing Unit}}
\newacronym{RAM}{RAM}{\textit{Random Access Memory}}
\newacronym{RSS}{RSS}{\textit{Resident Set Size}}
\newacronym{E/S}{E/S}{\textit{Entrada/Salida}}
\newacronym{AI}{AI}{\textit{Artificial Intelligence}}
\newacronym{CSV}{CSV}{\textit{Comma-Separated Values}}
\newacronym{PCAPNG}{PCAPNG}{\textit{Packet Capture Next Generation}}
\newglossaryentry{jitter}{name={jitter}, description={Variación temporal del retardo entre mediciones, paquetes o transacciones consecutivas}}
\newglossaryentry{entropia}{name={entropía}, description={Medida propuesta por Claude Shannon que cuantifica la incertidumbre promedio de una fuente discreta; en este trabajo se expresa en bits/byte para evaluar la variabilidad del contenido observable}}
\newglossaryentry{lineaBase}{name={Línea Base}, description={Escenario de referencia previo a la inserción del nodo intermedio y a la activación de MACsec}}
\newglossaryentry{puenteTransparente}{name={Puente transparente}, description={Escenario con el nodo SBC operando como puente de capa 2 sin protección MACsec activa}}

% =========================================================
% Estándares y organismos
% =========================================================
\newacronym{IEC}{IEC}{\textit{International Electrotechnical Commission}}
\newacronym{NIST}{NIST}{\textit{National Institute of Standards and Technology}}
\newacronym{DTIC}{DTIC}{\textit{Dirección de Tecnologías de Información y Comunicación}}
\newacronym{CIS}{CIS}{\textit{Center for Internet Security}}
\newacronym{GLPI}{GLPI}{\textit{Gestionnaire Libre de Parc Informatique}}
\newacronym{TI}{TI}{Tecnologías de Información}
\newacronym{TO}{TO}{Tecnologías de Operación}
